\appendix
\section{Axiom-by-axiom Satisfaction and Logical Independence}
\label{app:axiom-by-axiom}

This appendix upgrades the countermodel argument to an explicit model
satisfaction proof.  Let $T_{\rm PR}$ be the seven-postulate theory in PR-I
and let $L_{\rm PR}$ contain precisely the symbols fixed by those postulates.
Fix $R_A>0$ and $\Lambda_X^2>0$ and take the common reduct
\begin{align}
 \Psi&:\mathbb R^{1,3}\times(\mathbb R_w\times S^1_\theta)\to\mathbb C,\\
 ds_{\rm int}^2&=dw^2+R_A^2d\theta^2,\\
 \mathcal H_P&=L^2(\mathbb R_w\times S^1_\theta,R_A\,dw\,d\theta),\\
 O_{\rm int}&=O_X+O_\theta,\\
 O_X&=-\frac{d^2}{dw^2}+\Lambda_X^2\left(1+w^2+\frac34w^4\right),\\
 O_\theta&=-R_A^{-2}\frac{d^2}{d\theta^2}.
\end{align}
The usual Schrödinger and periodic domains make the two operators self-adjoint \cite{ReedSimon_1975}.  The quartic confinement gives a simple discrete radial spectrum. Define three expansions of this identical reduct:
\begin{equation}
 \mathcal M_{\rm pub}=(-1,\Pi_{13}),\qquad
 \mathcal M_{0}=(0,\Pi_{13}),\qquad
 \mathcal M_{15}=(-1,\Pi_{15}),
\end{equation}
where the first entry is $\sigma_8$ in
\begin{equation}
 N_\sigma=1+T_A^3+T_A^4+T_X(T_A^5+\sigma_8T_A^8).
\end{equation}

\begin{table}[H]
\centering
\small
\caption{Satisfaction of every frozen PR-I postulate.}
\label{tab:axiom-satisfaction}
\begin{tabularx}{\textwidth}{@{}p{0.08\textwidth}c c c X@{}}
\toprule
Axiom & $\mathcal M_{\rm pub}$ & $\mathcal M_0$ & $\mathcal M_{15}$ & Common satisfaction witness \\
\midrule
P1 & Pass & Pass & Pass & One identical master field and domain/codomain. \\
P2 & Pass & Pass & Pass & Identical $\mathbb R_w\times S^1_\theta$ topology. \\
P3 & Pass & Pass & Pass & Identical stationary metric, measure, and Laplacian. \\
P4 & Pass & Pass & Pass & Identical projection Hilbert space. \\
P5 & Pass & Pass & Pass & Identical separable internal operator. \\
P6 & Pass & Pass & Pass & Identical even quartic with $a_2=1$, $a_4=3/4$. \\
P7 & Pass & Pass & Pass & All five projection sectors exist; P7 does not fix $\sigma_8$ or the radial channel pair. \\
\bottomrule
\end{tabularx}
\end{table}

The first inequivalence witness is exact:
\begin{equation}
 c_{-1}=\frac{3354902985}{4211081216},\qquad
 c_0=\frac{52420359}{65798144},\qquad
 c_{-1}-c_0=\frac{9}{4211081216}\ne0.
\end{equation}
Both are dimensionless, so unit conventions cannot remove the difference.
For the spectral slices, the compressed radial traces are
$\lambda_1+\lambda_3$ and $\lambda_1+\lambda_5$.  Simplicity of the radial
spectrum gives $\lambda_3\ne\lambda_5$, and compressed trace is invariant
under a basis change.  Hence $\Pi_{13}$ and $\Pi_{15}$ are not quotient
equivalences.

\begin{theorem}[Logical independence]
The terminal orientation coefficient and the $\Pi_{13}$ channel selection are
not uniquely definable by P1--P7.
\end{theorem}
\begin{proof}
The three structures above have the same reduct to $L_{\rm PR}$ and therefore
satisfy the same P1--P7 sentences.  They assign different exact values to the
terminal cofactor or different nondegenerate spectral ranges to the boundary
projector.  If either selector were uniquely definable in $T_{\rm PR}$, every
expansion of the common reduct would have to assign it the same value, contrary
to these explicit expansions.
\end{proof}

Thus an axiom-by-axiom check strengthens rather than reverses the no-go result:
the positive singleton theorem is logically unavailable from the frozen
postulates.
